\section{The three-coloured flag surface}\label{sec:flag-surface}

The three views can be coupled more tightly than by their separate cycle
counts.  Fix $r\ne r'$ and a column $c$, and put
\[
 s=L(r,c),\qquad d=r_{r'}^{-1}(s),\qquad u=L(r,d).
\]
The three cells
\[
 (r,c,s),\qquad (r',d,s),\qquad (r,d,u)
\]
pairwise share a symbol, a column, and a row.  We call this triple a
\emph{flag}.  It determines one induced-cycle atom in each view: the row
atom on $\{r,r'\}$ through $c,d$, the column atom on $\{c,d\}$ through
$r,r'$, and the symbol atom on $\{s,u\}$ through $c,d$.

Retain the shared coordinate as an edge label: a column--symbol edge is
labelled by its row, a row--symbol edge by its column, and a row--column edge
by its symbol.  This distinction is essential when an intercalate produces
parallel edges.

\begin{theorem}[Flag surface and Gram identities]
\label{thm:flag-surface}
The $n^2(n-1)$ flags of a Latin square form the triangular faces of a closed,
possibly disconnected, three-coloured delta-complex $K(L)$.  A cycle atom of
length $\ell$ has link $C_{2\ell}$.  If $C_{V,\ell}$ is the number of such
atoms in view $V$, then
\[
 \chi(K(L))=
 \sum_{V\in\{R,C,S\}}\sum_{\ell=2}^n C_{V,\ell}
 -\frac{n^2(n-1)}2.
\]
Every intercalate is an isolated projective-plane delta-component.  If $L$
is FFF, every other nonempty component has negative Euler characteristic.

Let $B$ be the flag-by-cell incidence matrix, $P$ the flag-by-atom incidence
matrix, $Q$ the cell-by-atom support matrix, and $A$ the adjacency matrix of
the Latin-square graph on cells.  Then
\[
 B^{\mathsf T}B=3(n-1)I+2A,
 \qquad B^{\mathsf T}P=3Q.
\]
For $n>3$, $B^{\mathsf T}B$ is invertible and
\[
 P^{\mathsf T}P-
 9Q^{\mathsf T}(B^{\mathsf T}B)^{-1}Q\succeq0.
\]
\end{theorem}

\begin{proof}
Latinness gives $c\ne d$ and $s\ne u$, and the coloured incidences recover
$r,r',c$, so the stated flags are distinct.  A cycle atom of length $\ell$
is reached from its two line orientations and its $\ell$ cycle points, hence
has $2\ell$ incident flags.

A labelled pair of atoms in different views specifies one repeated diagonal
of a two-by-two rectangle.  Its two orientations are exactly the two incident
flags.  Thus every labelled edge lies in two faces.  Around an atom, the two
edge types that preserve it alternate along the induced cycle, giving the
link $C_{2\ell}$.  Hence $K(L)$ is closed.  With
$F=n^2(n-1)$ and $E=3F/2$, the Euler formula follows.  A $2$-cycle in any
view is one intercalate; its three atoms, six labelled edges, and four faces
form an isolated nonorientable component of Euler characteristic one, hence
a projective plane.  In an FFF residual component all vertex degrees are at
least eight.  Since $3F=\sum_v\deg(v)$ and $E=3F/2$, one obtains
$\chi=V-F/2\leq -V/3<0$.

Every cell lies in $3(n-1)$ flags.  Two different cells occur together in
two flags exactly when they share one row, column, or symbol, proving the
first matrix identity.  A cell in an atom support occurs in exactly three
flags incident with that atom, proving the second.  The Latin-square graph
has eigenvalues $3(n-1)$, $n-3$, and $-3$, so $B^{\mathsf T}B$ has
eigenvalues $9(n-1)$, $5n-9$, and $3(n-3)$ and is invertible for $n>3$.
The final assertion is the Schur complement of the positive-semidefinite
Gram matrix $[B\ P]^{\mathsf T}[B\ P]$.
\end{proof}

Let $\alpha_R,\alpha_C,\alpha_S$ be the fixed-point-free involutions on
flags that cross the labelled edge opposite the indicated atom colour.

\begin{proposition}[Dual-word traces]\label{prop:flag-traces}
For every $k\geq1$,
\[
 \operatorname{Fix}((\alpha_C\alpha_S)^k)
 =2\sum_{\ell\mid k}\ell C_{R,\ell},
\]
and cyclically in the other views.  If
$\beta=\alpha_R\alpha_C\alpha_S$, then
\[
 \operatorname{Fix}(\beta)=4I(L),
\]
where $I(L)$ is the intercalate count.  Every
$\operatorname{Fix}(\beta^k)$ is a main-class invariant.
\end{proposition}

\begin{proof}
Around a row atom of length $\ell$, the first product has two orbits of
length $\ell$, giving the divisor sum.  A flag is fixed by $\beta$ precisely
when crossing the three edge colours closes its two-by-two rectangle, which
means that the second diagonal also repeats a symbol.  This is an
intercalate, and each intercalate contributes its four flags.  Isotopy merely
relabels flags.  A permutation of the three coordinate colours sends
$\beta$ to a cyclic rotation of $\beta$ or $\beta^{-1}$; these have the same
fixed-point counts in every power.
\end{proof}

\begin{theorem}[Exact atom projection]\label{thm:flag-atom-projection}
Let $H=P^{\mathsf T}P$, let $v$ be the number of cycle atoms, and let $c$
be the number of connected components of $K(L)$.  Then
\[
 \operatorname{rank}_{\mathbb Q}H=v-2c.
\]
Let $N$ be the cell-by-coordinate incidence matrix, put
$U=Q^{\mathsf T}Q$, $W=Q^{\mathsf T}NN^{\mathsf T}Q$, and let
$\boldsymbol\ell$ be the vector of atom lengths.
Assume $n>3$ for the following Schur-complement assertions.  For
\[
 D=(n-1)(n-3)(5n-9),
\]
the scaled Schur complement has the integral form
\[
 DS=DH+6(n-1)W-3(n-1)(5n-9)U
       -32\boldsymbol\ell\boldsymbol\ell^{\mathsf T}\succeq0.
\]
Its nullity is at least $2c+1$.
\end{theorem}

\begin{proof}
If $Pz=0$, subtracting the equations of two triangular faces sharing an edge
of colours, say, $B$ and $C$ equates the values at their two opposite
$A$-vertices.  Along an arbitrary dual path, crossing an edge of colour $A$
therefore replaces the $A$-vertex by one with the same value, whereas crossing
an edge of colour $B$ or $C$ leaves the $A$-vertex literally unchanged because
it belongs to the shared primal edge.  Dual connectivity consequently makes
$z$ constant on all $A$-vertices in a component.  The same argument applies
cyclically to $B$ and $C$.  Each face equation says that the three resulting
constants sum to zero.  Thus the kernel has exactly two dimensions per
component, which proves the rank formula.

The cell-coordinate matrix satisfies $NN^{\mathsf T}=3I+A$.  Moreover
\[
 (B^{\mathsf T}B)^{-1}
 =\frac5{3(5n-9)}I
 -\frac2{3(n-3)(5n-9)}A
 +\frac8{9(n-1)(n-3)(5n-9)}J.
\]
Substitution in $S=H-9Q^{\mathsf T}(B^{\mathsf T}B)^{-1}Q$ gives the
displayed integer formula.  Positive semidefiniteness follows because this
is $P^{\mathsf T}$ times the orthogonal projector away from
$\operatorname{col}(B)$ times $P$.  The componentwise colour differences
give $2c$ kernel vectors.  One further vector is constant on all atoms of
one colour: its image under $P$ is the all-one flag vector, which lies in
$\operatorname{col}(B)$.
\end{proof}

\begin{corollary}[Topological and sign-cut consequences]
\label{cor:flag-topological-sign-cut}
On an orientable flag component every cycle of every odd-length word in the
dual involutions has even length.  If $n$ is even and $L$ is FFF, then
\[
 \chi(K(L))\equiv \frac n2\pmod2.
\]
For even $n$, every F-view has at least
\[
 n(n-1)-\left\lfloor\frac{n^2}{4}\right\rfloor
\]
cycle atoms.  Thus a hypothetical order-$10$ FFF square has at least $65$
atoms in each view and $195$ in total.
\end{corollary}

\begin{proof}
Coherent face signs change across every edge, so an odd word reverses sign
and can return only after an even number of iterations.  For the Euler
parity, all atoms are even cycles and the global sign relation of
Proposition~\ref{prop:aggregate-profile-baseline} gives
$v\equiv n(n-1)/2\pmod2$; moreover $n^2(n-1)/2$ is even, so the face
contribution to Euler characteristic is even.  Finally, by
Proposition~\ref{prop:line-sign-cut}, line pairs with an
odd number of cycles form a cut of size at most $\lfloor n^2/4\rfloor$.
Such a pair contributes at least one atom and every other pair at least two.
\end{proof}

The aggregate consequences in this corollary are not sufficient for an
order-$10$ obstruction.  An exact numerical repository witness uses the
disjoint-union summaries of frozen order-$8$ FFF flag surfaces 1 and 2 and
one isolated projective-plane intercalate component.  It has $900$ flags, $591$ atoms,
$176$ components and Euler characteristic $141$;
Theorem~\ref{thm:flag-atom-projection} then forces
$\operatorname{rank}H=239$.  In each colour its $169$ length-two and $28$
length-four atoms can be regrouped into $45$ formal degree-$10$ line-pair
partitions, with $25$ odd-cycle-count pairs forming the size of a $5|5$ cut.
The three cut parameters are $(5,5,5)$ and therefore have the required odd
sum.  This is a decoupled aggregate consistency witness only.  The formal
line-pair labels are not shown compatible with the flag faces, labelled
edges, or $H$ incidences; nor does it supply a common set of $100$ Latin
cells, $Q$, coordinate-multiplicity, or Schur data.  It shows precisely that
the listed aggregate and componentwise conditions cannot by themselves
decide order $10$ by itself.  Pair-label/flag-incidence compatibility is an
intermediate obstruction layer not excluded by this construction.

\begin{theorem}[Labelled line-pair Gram reduction]
\label{thm:labelled-pair-gram}
Assume $n>3$, and put $m=\binom n2$.  For a flag $\phi$, let $\lambda_V(\phi)$ be its
unordered line-pair label in colour $V$, and let $Z_V$ be the corresponding
flag-by-pair incidence matrix.  Then
\[
 Z_V^{\mathsf T}Z_V=2nI_m,
 \qquad W^{VW}:=\tfrac12 Z_V^{\mathsf T}Z_W\in\{0,1,2\}^{m\times m}.
\]
Every row and column of $W^{VW}$ is the weighted edge vector of a two-factor
of $K_n$, and its component lengths are the cycle lengths of the associated
two-line permutation.  If $E$ is the unsigned vertex-edge incidence matrix
of $K_n$, then
\[
 EW^{VW}=2J,\qquad W^{VW}E^{\mathsf T}=2J.
\]
Moreover
\[
 K=\begin{bmatrix}
 nI&W^{RC}&W^{RS}\\
 (W^{RC})^{\mathsf T}&nI&W^{CS}\\
 (W^{RS})^{\mathsf T}&(W^{CS})^{\mathsf T}&nI
 \end{bmatrix}=\tfrac12 Z^{\mathsf T}Z\succeq0.
\]
Writing $C$ for the restriction of this block matrix to three copies of
$\ker E$, one has
\[
 \operatorname{rank}_{\mathbb Q}K=3n-2+
 \operatorname{rank}_{\mathbb Q}C.
\]
The Schur projection from Theorem~\ref{thm:flag-atom-projection}, obtained by
removing the coarse flag-by-cell
incidence satisfies
\[
 \Sigma\succeq0\Longleftrightarrow K\succeq0,
 \qquad \operatorname{rank}_{\mathbb Q}\Sigma=
 \operatorname{rank}_{\mathbb Q}K-1.
\]
Finally,
\[
 t(L)=\operatorname{tr}\!\left(
 W^{RC}W^{CS}(W^{RS})^{\mathsf T}\right)
\]
is a main-class invariant.
\end{theorem}

\begin{proof}
Each line-pair label occurs in $2n$ flags.  A cross-codegree counts repeated
diagonals in the rectangle selected by two pair labels; its value is zero,
one, or two.  Equivalently, $W^{RC}$ is the sum over symbols of the
permutation matrices induced on unordered row pairs by the corresponding
row--column matchings, and cyclically for the other colours.  This proves the
two-factor statement and the equations with $E$.

The displayed block matrix is its defining Gram matrix.  Since
$EE^{\mathsf T}=(n-2)I+J$, one has the orthogonal decomposition
\[
 \mathbb Q^m=\langle\mathbf1\rangle\oplus
 E^{\mathsf T}(\mathbf1^\perp)\oplus\ker E.
\]
Every $W$ acts by $n$ on constants, annihilates endpoint contrasts, and
preserves $\ker E$.  The constant block has rank one, the three endpoint
blocks have total rank $3(n-1)$, and the rank formula follows.

For the Schur statement, direct flag counting gives
$B^{\mathsf T}Z_V=3N_VE$, where $N_V$ records the $V$-line of a cell.
Substitution of the closed inverse from
Theorem~\ref{thm:flag-atom-projection} diagonalizes $\Sigma$ on the same
decomposition: it is zero on the common constant direction, positive
definite on the three endpoint-contrast spaces, and equals $2C$ on the
cycle spaces.  This proves both assertions.  Independent row, column, and
symbol relabellings conjugate the three pair-label sets, while parastrophy
only permutes the colours and transposes blocks.  The symmetric triangle
contraction $t(L)$ is therefore invariant.
\end{proof}

This theorem isolates the full nontrivial common-Gram information in
dimension $3(m-n)$, but it does not characterize realizable Latin squares.
In particular, the coarse cell projection cannot supply a stronger
pair-matrix obstruction; finer common cell support or another nonlinear
compatibility is required.

\begin{proposition}[Full pair-label tensor]
\label{prop:pair-label-tensor}
For unordered row-, column-, and symbol-pair labels $e,f,g$, let
\[
 T(e,f,g)=\#\{\phi:\lambda_R(\phi)=e,
 \lambda_C(\phi)=f,\lambda_S(\phi)=g\}.
\]
Every entry of $T$ lies in $\{0,1,4\}$.  Its two-dimensional marginals are
$2W^{RC},2W^{RS},2W^{CS}$.  More precisely, a marginal entry zero has no
completion, an entry two has two distinct completions of multiplicity one,
and an entry four has one completion of multiplicity four; the latter is
exactly an intercalate.

The multiplicity-one support is a closed three-coloured pair-label
pseudocomplex $X_\lambda(L)$.  At a line-pair vertex its link is the disjoint
union of the cycles $C_{2\ell}$ belonging to the nonintercalate induced
cycles of lengths $\ell$ on that line pair.  Consequently $L$ is FFF if and
only if every component of every vertex link of $X_\lambda(L)$ has length
divisible by four.

For a fixed row $r$, contracting $T$ over all row-pair labels incident with
$r$ and thresholding entries at two gives exactly the two-subset action of
the row permutation $c\mapsto L(r,c)$.  Cyclic contractions recover the
column and symbol line maps.  Since the two-subset action of $S_n$ is
faithful for $n\geq3$, the fixed-coordinate map $L\mapsto T(L)$ is injective
for $n\geq3$.  This restriction is necessary: the two coordinate-labelled
order-$2$ squares have the same single tensor entry, of value $4$.
\end{proposition}

\begin{proof}
Fix a row pair and a column pair.  A flag exists exactly when a diagonal of
the selected rectangle repeats a symbol.  If precisely one diagonal repeats,
its two orientations give two flags with different symbol-pair labels;
equality of those labels would make the other diagonal repeat as well.  If
both diagonals repeat, the rectangle is an intercalate and all four flags
have the same three pair labels.  This proves the local law and every
marginal statement, cyclically in the colours.

After deleting multiplicity-four faces, each active two-label edge has the
two multiplicity-one completions just described, so it lies in exactly two
faces.  Around a fixed pair label, the flag-surface links of distinct nonintercalate
cycle atoms remain disjoint and are $C_{2\ell}$.  Splitting a singular
pair-label vertex by these link components recovers the nonintercalate part
of the flag surface in Theorem~\ref{thm:flag-surface}.  Finally, an induced
cycle length $\ell>2$ is even
exactly when $2\ell$ is divisible by four; the omitted intercalate cycles
have length two and are already F-compatible.

For the contraction claim, fix a column pair $\{c,d\}$.  The pair
$\{L(r,c),L(r,d)\}$ is the unique symbol-pair label of contracted
multiplicity at least two: the positive multiplicities are $(2,1,1)$ unless
the two diagonal completions use the same partner row, in which case the
rectangle is an intercalate and the sole multiplicity is four.  Thus the
threshold matrix is the edge lift of the row permutation.  For distinct
$d_1,d_2\ne c$, the intersection of the two image edges of
$\{c,d_1\}$ and $\{c,d_2\}$ is the singleton $\{L(r,c)\}$, proving
faithfulness and reconstruction.  The cyclic cases are identical.
\end{proof}

\begin{theorem}[Smith torsion and total unimodularity]
\label{thm:pair-tensor-smith}
Fix one pair-label coordinate of the tensor $T$ in
Proposition~\ref{prop:pair-label-tensor}, and let $T_e$ be the resulting
integral matrix.  Delete its identically zero rows and columns and denote the
active square matrix by $\widehat T_e$.  If the associated two-line
permutation has $t$ cycles of
length two, $o$ odd cycles, and $q$ even cycles of length at least four, then,
up to unit Smith factors,
\[
 \operatorname{coker}(\widehat T_e)\cong
 \mathbb Z^q\oplus(\mathbb Z/2\mathbb Z)^o
 \oplus(\mathbb Z/4\mathbb Z)^t.
\]
In particular,
\[
 o=\operatorname{rank}_{\mathbb Q}(T_e)
   -\operatorname{rank}_{\mathbb F_2}(T_e\bmod2)
   -\#\{(f,g):T_e(f,g)=4\}.
\]
Consequently a Latin square is FFF if and only if this expression is zero for
every slice in all three tensor directions.

The ranks in the displayed formula are unchanged when zero rows and columns
are deleted; the full cokernel merely has the corresponding additional free
factors.

Equivalently, define $N_e(f,g)=1$ when $T_e(f,g)>0$ and zero otherwise.  A
Latin square is FFF if and only if every $N_e$ in all three directions is
totally unimodular; equivalently, every such matrix is balanced.  We use the
standard terminology of Camion and Berge~\cite{Camion1965,Berge1972}; the
equivalence asserted here is for these special two-regular slice blocks, not
for arbitrary zero-one matrices.
\end{theorem}

\begin{proof}
The slice normal form from Proposition~\ref{prop:pair-label-tensor} is a
direct sum of $[4]$ over the length-two cycles and
$I_\ell+P_\ell$ over all longer cycles.  The cokernel of the latter block has
generators $x_i$ with relations $x_i+x_{i+1}=0$.  Hence every generator is
$(-1)^ix_0$, and closing the cycle leaves a free generator when $\ell$ is
even and imposes $2x_0=0$ when $\ell$ is odd.  This proves the Smith and rank
formulas; the entries equal to four are precisely the intercalates.

After normalization, the blocks are $[1]$ and $I_\ell+P_\ell$.  For odd
$\ell$ the latter has determinant $2$ and is not totally unimodular.  For
even $\ell$ it is the unsigned vertex-edge incidence matrix of a bipartite
cycle; changing signs on one vertex class makes it an oriented incidence
matrix, which is totally unimodular.  Direct sums preserve this property.
The balanced-matrix statement is the same forbidden odd-cycle condition for
these two-regular incidence blocks.
\end{proof}

\begin{theorem}[Triple endpoint contraction]
\label{thm:triple-endpoint-parity}
Let $E_R,E_C,E_S$ be the unsigned vertex--edge incidence matrices of the
complete graphs on the row, column, and symbol sets, respectively, and put
\[
 D=(E_R\otimes E_C\otimes E_S)T,
 \qquad H=D\bmod 2.
\]
Every one-dimensional line of $D$ has sum $8(n-1)$.  If
$x=(r,c,L(r,c))$ is a Latin cell and $\iota(x)$ is the number of
intercalates containing it, then
\[
 D(r,c,L(r,c))=3(n-1)+\iota(x),
 \qquad
 \sum_{s\ne L(r,c)}D(r,c,s)=5(n-1)-\iota(x).
\]
Consequently $H$ lies in the tensor product of the three even-sum point
spaces over $\mathbb F_2$, each of its mode ranks is at most $n-1$, and its
values on the Latin incidence positions recover the cellwise intercalate
parity after the table is reconstructed by
Theorem~\ref{thm:endpoint-characterization}.

More precisely, write $T=X+4Y$ with $X=\mathbf1_{T=1}$ and
$Y=\mathbf1_{T=4}$.  If
$A=(E_R\otimes E_C\otimes E_S)X$ and
$B=(E_R\otimes E_C\otimes E_S)Y$, then on every Latin cell $x$,
\[
 B(x)=\iota(x),\qquad A(x)=3\bigl(n-1-\iota(x)\bigr).
\]
In particular $0\leq\iota(x)\leq n-1$.
\end{theorem}

\begin{proof}
Fix row and column endpoints $(r,c)$.  Summing over the symbol endpoint
contributes two endpoints for every symbol edge.
Proposition~\ref{prop:pair-label-tensor} and
Theorem~\ref{thm:labelled-pair-gram} therefore give
\[
 \sum_sD(r,c,s)
 =4\sum_{e\ni r,\,f\ni c}W^{RC}(e,f)
 =8(n-1).
\]
The other two directions follow cyclically.  Reduction modulo two proves the
even-line and rank statements.

Every cell lies in exactly $3(n-1)$ flag faces.  A nonintercalate flag has
exactly its three face cells in its endpoint box, so it contributes at $x$
exactly when the face contains $x$.  The four flags of an intercalate have
one common triple of pair labels.  All four endpoint boxes contain each of
the four intercalate cells, whereas only three of the four faces contain a
fixed cell.  Hence every intercalate through $x$ contributes one extra unit,
which proves the first displayed cell formula.  Subtracting it from the line
sum proves the complementary formula.

Finally, every multiplicity-four entry of $T$ is one intercalate and its
endpoint box contains precisely its four Latin cells.  Thus
$B(x)=\iota(x)$; the formula for $A$ follows from $D=A+4B$.  Its
nonnegativity gives the local upper bound.
\end{proof}

\begin{theorem}[Colour-parity holonomy]
\label{thm:colour-parity-holonomy}
Let $K^*(L)$ be the dual cellulation of the flag surface.  Assign to every
dual edge of colour $V\in\{R,C,S\}$ the corresponding standard basis vector
$e_V\in\mathbb F_2^3$, and denote the resulting one-cochain by $\omega$.
Then $\omega$ is a cocycle if and only if $L$ is FFF.

More precisely, the coordinate cochain $\omega_R$ is closed exactly when the
column and symbol views are F, and cyclically.  On every connected FFF
component,
\[
 [\omega_R]+[\omega_C]+[\omega_S]=w_1(K(L)).
\]
If the three colour classes span a subspace of dimension $h$ in
$H^1(K(L);\mathbb F_2)$, then the number of flags in that component is
divisible by $2^{3-h}$.

On every intercalate component, which is a projective plane, the three colour
classes are equal to the same nonzero class.  Finally, if $L$ is FFF and
$n\equiv2\pmod4$, at least one component has nontrivial colour-parity
holonomy.
\end{theorem}

\begin{proof}
An atom in the row view with induced cycle length $\ell$ has a dual boundary
containing $\ell$ column-coloured and $\ell$ symbol-coloured edges.  Therefore
\[
 (\delta\omega)(a)=\ell(e_C+e_S),
\]
and the two cyclic analogues hold in the other views.  This vanishes exactly
when the atom length is even.  Thus the three colour-stratified coboundaries
locate every odd-cycle atom and retain the full three-view pattern.  By
contrast, the three booleans recording whether the coordinate cochains are
globally closed are only the pairwise conjunctions of the F-view bits and do
not distinguish patterns with at least two failing views.

Orient every primal flag triangle first by the fixed colour order $(R,C,S)$.
The two induced orientations on a shared labelled edge agree, so the
transition cochain is one on every dual edge and represents $w_1$.  Since
$\omega_R+\omega_C+\omega_S$ is precisely this constant-one cochain, the
displayed cohomology identity follows.

Let $H\leq\mathbb F_2^3$ be the holonomy subgroup obtained by integrating
$\omega$ around closed dual paths.  Its dimension is $h$.  Path integration
labels flags by $\mathbb F_2^3/H$.  A word in the three dual edge involutions
acts bijectively on the component and translates every label by the same
element.  Because the three basis vectors generate the quotient, all
$2^{3-h}$ label fibres have equal size, proving divisibility.

The dual graph of an intercalate component is $K_4$, with its three colour
classes as the three perfect matchings.  Every triangle uses all three
colours, so each coordinate evaluates nontrivially on the unique nonzero
mod-two homology class; all three equal $w_1\ne0$.

For the final statement, Corollary~\ref{cor:flag-topological-sign-cut} gives
$\chi(K(L))\equiv n/2\pmod2$ for even-order FFF squares.  If
$n\equiv2\pmod4$, the total Euler characteristic is odd.  Since every closed
orientable surface has even Euler characteristic, some component is
nonorientable.  The displayed identity then forces nonzero colour-class
span on that component.
\end{proof}

\begin{definition}[Exact masks and compatible potentials]
\label{def:exact-mask-potentials}
Let $F$ be a connected component of the dual flag cellulation.  For
$m=(m_R,m_C,m_S)\in\mathbb F_2^3$, put
\[
 \omega_m=m_R\omega_R+m_C\omega_C+m_S\omega_S.
\]
The mask $m$ is \emph{exact on $F$} if there is a binary zero-cochain $p_m$
on the flags of $F$ with $\delta p_m=\omega_m$, equivalently
\[
 p_m(\alpha_Vf)=p_m(f)+m_V
 \qquad(V\in\{R,C,S\}).
\]
The exact masks form a vector space $M_F$.  A \emph{compatible family of
potentials} is a linear choice $m\mapsto p_m$ for $m\in M_F$, modulo
componentwise constants.  Such a family is obtained by choosing potentials
for a basis of $M_F$, normalizing them at one base flag, and extending
linearly.
\end{definition}

Three-digit masks are written in tuple order $(R,C,S)$.  When integer mask
codes are used below, the code of $m$ is $m_R+2m_C+4m_S$; hence codes
$3,5,6$ correspond respectively to tuples $110,101,011$.  In particular,
an integer's usual binary display must not be read as an $(R,C,S)$ tuple.

The componentwise multiset of flag count, Euler characteristic,
orientability, and exact linear relations among the three colour classes is
preserved by isotopy.  Parastrophy applies one global permutation of the
three colours.  Sorting these tuples and minimizing over the six global
colour permutations therefore gives a main-class invariant.

\begin{computationaltheorem}[Order-$8$ holonomy profiles]
\label{cthm:holonomy-profile}
On the $230$ frozen order-$8$ FFF main-class representatives, the canonical
component-holonomy profile takes exactly $152$ values.  Exactly $109$ values
are singleton fibres; the maximum fibre size is five.  The collision
histogram by fibre size is
\[
 1:109,\qquad 2:22,\qquad 3:12,\qquad 4:4,\qquad 5:5.
\]
\end{computationaltheorem}

This invariant records the global gluing of the alternating bipartitions
described by Theorem~\ref{thm:colour-parity-holonomy}.  The collisions in the computational
theorem show that it is not a complete classifier, and no injectivity in
higher orders is asserted.

\begin{theorem}[Oriented holonomy-to-cell lift]
\label{thm:holonomy-cell-lift}
Let $I$ be an intercalate with rows $r,r'$, columns $c,c'$, and symbols
$s,t$, displayed as
\[
 \begin{array}{c|cc}
      &c&c'\\ \hline
 r    &s&t\\
 r'   &t&s
 \end{array}.
\]
On its four-flag component, the exact colour masks are
$000,011,101,110$.  For each nonzero exact mask $m$, choose a binary
potential on the flags, replace its values by signs, and push those signs to
the four Latin cells by summing over incident flags.  Up to an independent
global sign in each mask, the resulting cell vectors are
\[
 u_{011}=(1,1,-1,-1),\qquad
 u_{101}=(-1,1,-1,1),\qquad
 u_{110}=(1,-1,-1,1).
\]
Thus $u_{011}$ has zero column and symbol marginals, $u_{101}$ has zero row and
symbol marginals, and $u_{110}$ has zero row and column marginals.  The
remaining marginal is in each case twice an oriented endpoint difference,
and
\[
 u_{011}u_{101}u_{110}=-\mathbf 1_I
\]
pointwise after the displayed gauge choice.

If these vectors are assembled as cell-by-intercalate matrices, their three
entrywise absolute values coincide with the intercalate incidence matrix
$J$.  Consequently
\[
 J\mathbf 1(x)=\iota(x)
 =\operatorname{endpoint}(Y)(x)
\]
on every Latin cell $x$, where $Y$ is the multiplicity-four layer in
Theorem~\ref{thm:triple-endpoint-parity}.
\end{theorem}

\begin{proof}
Every flag of the intercalate contains exactly three of its four cells, and
each cell $x$ is omitted by a unique flag $f_x$.  A potential for a nonzero
exact mask has two positive and two negative signs: its two selected colour
matchings form a four-cycle in the dual $K_4$, and the potential alternates
around that cycle.  Their total is zero, so
the pushforward at $x$ equals the negative of the sign at $f_x$.  It
therefore has absolute value one on all four cells.  Choosing the three
potential gauges gives the displayed vectors, from which all marginal and
pointwise-product identities follow directly.

The common absolute-value column is precisely the incidence vector of $I$.
Summing these columns counts the intercalates through $x$.  The final equality
is the multiplicity-four endpoint formula of
Theorem~\ref{thm:triple-endpoint-parity}.
\end{proof}

The theorem supplies an oriented lift of the complete multiplicity-four
attachment layer.  It imposes no condition on the multiplicity-one layer
$X$ and hence is not, by itself, an order-$10$ obstruction.

\begin{theorem}[Exact-mask cell pushforwards]
\label{thm:exact-mask-cell-pushforward}
Let $F$ be any connected flag-surface component, let
$m=(m_R,m_C,m_S)$ be an exact colour mask on $F$, and choose a binary
potential $p_m$ on its flags.  Define
\[
 u_m(x)=\sum_{\substack{f\in F\\x\in f}}(-1)^{p_m(f)}
\]
on Latin cells.  If $m_V=1$, then every line marginal of $u_m$ in view $V$
is zero.  Thus an exact mask kills all marginal families indexed by its set
bits.  Complementing the potential changes $u_m$ only by a global sign.
\end{theorem}

\begin{proof}
Order the three cells in every flag by the defining flag rectangle.  For a
fixed view $V$ and cell role, there is a fixed-point-free involution preserving
the $V$-line of that role.  It is either the dual involution $\alpha_V$ or a
conjugate $\alpha_W\alpha_V\alpha_W$: row roles zero and two use $\alpha_R$
and row role one uses $\alpha_C\alpha_R\alpha_C$, with the cyclic analogues
for columns and symbols.  These statements follow immediately by writing the
three cells of the adjacent flag rectangles.

Across $\alpha_V$ the potential changes by $m_V$; across the displayed
conjugate it changes by $m_W+m_V+m_W=m_V$.  If $m_V=1$, the two signs in each
role pair are opposite while their selected cells lie on the same $V$-line.
Each role sum, and hence the full line marginal, vanishes.
\end{proof}

Unlike Theorem~\ref{thm:holonomy-cell-lift}, this statement applies to every
component and therefore reaches the multiplicity-one surface.  It remains a
componentwise linear consequence of exactness, not a joint $X/Y$
compatibility theorem.

\begin{theorem}[Cellwise holonomy Fourier transform]
\label{thm:cellwise-holonomy-fourier}
Let the exact colour-mask space of a connected flag component have basis
$m_1,\ldots,m_d$, and choose compatible binary potentials $p_i$.  Label a
flag by $q(f)=(p_1(f),\ldots,p_d(f))\in\mathbb F_2^d$, and let
$N_x(a)$ count the flags of label $a$ incident with a Latin cell $x$.  For
$c\in\mathbb F_2^d$, let $u_c$ be the cell pushforward of the potential
$\sum_i c_i p_i$.  Then
\[
 u_c(x)=\sum_{a\in\mathbb F_2^d}(-1)^{c\cdot a}N_x(a).
\]
Consequently,
\[
 N_x(a)=2^{-d}\sum_c(-1)^{c\cdot a}u_c(x),\qquad
 \sum_c u_c(x)^2=2^d\sum_a N_x(a)^2.
\]
In particular, every inverse signed sum is a nonnegative multiple of $2^d$.
\end{theorem}

\begin{proof}
Group the flags in the definition of $u_c(x)$ by their label $a$.  This gives
the first formula.  The other two are Fourier inversion and Parseval on the
elementary abelian group $\mathbb F_2^d$.  A basis change relabels the group,
while complementing a potential translates its labels and changes character
signs, so the asserted constraints are intrinsic.
\end{proof}

For an intercalate component, $d=2$ and a supported cell sees three of the
four flag labels.  Parseval becomes $3^2+1^2+1^2+1^2=4\cdot3$, recovering
the oriented intercalate vectors of
Theorem~\ref{thm:holonomy-cell-lift}.  In general the theorem couples all exact masks at a cell,
but it does not yet relate different connected components.

\begin{theorem}[Affine transport of cell roles]
\label{thm:cell-role-affine-transport}
In the notation of Theorem~\ref{thm:cellwise-holonomy-fourier}, let
$N_{x,j}(a)$ count label-$a$ flags in which $x$ has cell role $j$, and let
$b_V=((m_1)_V,\ldots,(m_d)_V)$.  Then
\[
 N_{x,2}(a+b_R)=N_{x,0}(a),\quad
 N_{x,2}(a+b_C)=N_{x,1}(a),\quad
 N_{x,1}(a+b_S)=N_{x,0}(a).
\]
Hence the three role totals are equal.  With
$\delta=b_R+b_C+b_S$, one has
$N_{x,0}(a+\delta)=N_{x,0}(a)$.  If $\delta\ne0$, the total component
incidence at $x$ is divisible by six.  This holds in particular for every
orientable component.
\end{theorem}

\begin{proof}
The dual involution $\alpha_R$ exchanges shared roles $(0,2)$ and adds
$b_R$ to the flag label.  Likewise, $\alpha_C$ exchanges $(1,2)$ and adds
$b_C$, while $\alpha_S$ exchanges $(0,1)$ and adds $b_S$.  This gives the
three displayed bijections.
Their two routes from role zero to role two yield invariance under $\delta$.
If $\delta\ne0$, its translation orbits have size two, so each role total is
even.  There are three equal role totals.  Finally, orientability makes mask
$111$ exact by Theorem~\ref{thm:colour-parity-holonomy}; therefore not every
exact mask has even weight and $\delta\ne0$.
\end{proof}

\begin{theorem}[Cross-component cell-parity cover]
\label{thm:cross-component-cell-parity}
For a flag-surface component $F$, let $t_F(x)$ be its common role mass from
Theorem~\ref{thm:cell-role-affine-transport} at
the Latin cell $x$, and put $p_F(x)=t_F(x)\bmod2$.  If
$e^R_F(r),e^C_F(c),e^S_F(s)$ count the correspondingly labelled dual edges,
then
\[
 \sum_{x\in r}t_F(x)=2e^R_F(r),\quad
 \sum_{x\in c}t_F(x)=2e^C_F(c),\quad
 \sum_{x\in s}t_F(x)=2e^S_F(s).
\]
Thus every $p_F$ has even sum on each row, column, and symbol class.  Moreover,
\[
 \sum_F t_F(x)=n-1,
 \qquad
 \sum_{F:\delta_F=0}p_F=(n-1)\mathbf1
 \quad\text{over }\mathbb F_2.
\]
For every colour $V$ and corresponding point label $a$, the component edge
counts also satisfy $\sum_F e^V_F(a)=n(n-1)/2$.
If $\bar\iota$ is the cellwise intercalate-degree vector modulo two, then
\[
 \sum_{\substack{F\ \mathrm{nonintercalate}\\\delta_F=0}}p_F
   =(n-1)\mathbf1+\bar\iota
   =\left.(E_R\otimes E_C\otimes E_S)X\right|_L\pmod2.
\]
Consequently, at even order every cell lies in an odd number of
$\delta_F=0$ components with odd role mass.  If its intercalate degree is
even, an odd number of these components are nonintercalate.
\end{theorem}

\begin{proof}
For a fixed row, count $t_F$ using role zero.  The corresponding faces are
exactly the faces incident with the $\alpha_R$ edges carrying that row label,
and every such edge has two incident faces.  This proves
$\sum_{x\in r}t_F(x)=2e^R_F(r)$.  Using role two and $\alpha_C$ gives the
column identity; using role zero and $\alpha_S$ gives the symbol identity.
Their reductions modulo two give the asserted even-line conditions.

Every cell occurs exactly $n-1$ times in each role globally, which gives the
first displayed equality.  Summing any one of the three line identities over
$F$ and then over the $n$ cells of the fixed line gives
$2\sum_F e^V_F(a)=n(n-1)$, which proves the global edge budget.
Theorem~\ref{thm:cell-role-affine-transport} makes $t_F(x)$ even when
$\delta_F\ne0$, giving the second displayed equality.  An intercalate
component has $\delta_F=0$ and contributes its four-cell incidence vector,
so their sum is $\bar\iota$.  Finally, the multiplicity-one formula in
Theorem~\ref{thm:triple-endpoint-parity} is
$((E_R\otimes E_C\otimes E_S)X)(x)=3(n-1-\iota(x))$ on Latin cells.  Reduction
modulo two proves the last identity and its parity consequences.
\end{proof}

\begin{theorem}[Gauge selection and cubic cell lift]
\label{thm:gauge-invariant-cell-lift}
Let $M_F\leq\mathbb F_2^3$ be the exact-mask space of a connected flag
component.  Choose mask potentials compatibly, and let $v_{m,j}(x)$ be the
signed pushforward at the Latin cell $x$ restricted to cell role $j$.  Then
\[
 v_{m,2}=(-1)^{m_R}v_{m,0},\qquad
 v_{m,2}=(-1)^{m_C}v_{m,1},\qquad
 v_{m,1}=(-1)^{m_S}v_{m,0}.
\]
Consequently every odd-weight exact mask has identically zero full cell
pushforward $u_m=\sum_jv_{m,j}$.

A compatible gauge change is a linear functional
$\lambda:M_F\to\mathbb F_2$ and sends
$u_m\mapsto(-1)^{\lambda(m)}u_m$.  The formal monomial
$\prod_m u_m^{e_m}$ has trivial transformation character exactly when
\[
 \sum_m(e_m\bmod2)m=0.
\]
For an evaluated monomial that is not identically zero, this is also
equivalent to gauge invariance.  An identically zero evaluated monomial is
invariant regardless of its formal charge.
In particular, independent component gauges admit no nonzero linear signed
cancellation across components without anchors.

If the full even-mask plane $\{0,3,5,6\}$ is exact, then
\[
 K_F(x)=-u_3(x)u_5(x)u_6(x)
\]
is gauge invariant.  With $t_F(x)$ the common role mass,
\[
 |K_F(x)|\leq t_F(x)^3,\qquad
 K_F(x)\equiv t_F(x)^3\pmod4,
\]
and $8\mid K_F(x)$ when $t_F(x)$ is even.  On every intercalate component,
$K_F=1$ on its four cells.
\end{theorem}

\begin{proof}
The three displayed role identities are the character transforms of the
three affine bijections in Theorem~\ref{thm:cell-role-affine-transport}.
Comparing the first with the other two reverses the sign of $v_{m,0}$ exactly
when $|m|$ is odd, forcing that integer and hence all role sums to vanish.

Two compatible potential families differ by a componentwise constant linear
functional $\lambda$.  A monomial therefore acquires the sign obtained by
pairing $\lambda$ with its displayed mask charge.  If that charge is nonzero,
some linear functional takes value one on it; a nonzero evaluated monomial
then changes sign.  This proves the character criterion and its qualified
value-level converse.  Since the gauges of distinct components are independent, a
charged linear contribution cannot be cancelled intrinsically by another
component.

For the three nonzero even masks the role formula gives
$u_3=v_{3,0}$, $u_5=-v_{5,0}$, and $u_6=v_{6,0}$.  Thus $K_F$ is the product
of the three nontrivial Walsh coefficients of the four-point role-zero label
distribution.  If its total is $t$, write these coefficients as $t-2s_i$.
Their absolute values are at most $t$.  For even $t$ all three are even.  For
odd $t$, every nonzero four-point label belongs to the negative half of
exactly two nontrivial characters, so $\sum_i s_i$ is even; expanding
$\prod_i(t-2s_i)$ modulo four proves the congruence.  The intercalate value is
the compatible-sign specialization of Theorem~\ref{thm:holonomy-cell-lift}.
\end{proof}

The cubic invariant contains genuine nonintercalate information and does not
sum merely to the intercalate degree.  For example, in the reduced FFF square
\[
\begin{pmatrix}
0&1&2&3\\1&3&0&2\\2&0&3&1\\3&2&1&0
\end{pmatrix}
\]
the cell $(0,0,0)$ has one nonintercalate contribution $8$ and one
intercalate contribution $1$, so $\sum_FK_F=9\ne1=\iota$.  Thus the
cross-component parity cover in
Theorem~\ref{thm:cross-component-cell-parity} remains
the canonical binary cover; unanchored linear signs and this naive cubic sum
do not provide a stronger gluing law.

\begin{definition}[Equivariant gauge-selection rule]
A \emph{gauge-selection rule} assigns, to every coordinate-labelled Latin
square and every connected flag component, one compatible potential family
from Definition~\ref{def:exact-mask-potentials}.  It is
\emph{isotopy-equivariant} if every colour-preserving isotopy transports each
selected potential to the selected potential on the image component.  A rule
constructed only from the complete endpoint-lift data is understood to use no
additional root flag, distinguished cell, or external ordering.
\end{definition}

\begin{theorem}[Quadratic Gram lift and equivariant gauge obstruction]
\label{thm:quadratic-gram-gauge-obstruction}
For an exact mask $m$ on a flag component $F$, put
\[
 Q_{F,m}=u_{F,m}u_{F,m}^{\mathsf T},\qquad
 G_m=\sum_{F:m\in M_F}Q_{F,m}.
\]
These integral cell kernels are independent of all compatible potential
gauges, symmetric, and positive semidefinite.  If the $V$ bit of $m$ is set,
every $V$-line indicator lies in their kernel.  In particular, for each
nonzero even mask $m\in\{3,5,6\}$,
\[
 \operatorname{rank}G_m\leq(n-1)^2,\qquad
 G_m(x,y)\equiv
 \sum_{F:m\in M_F}p_F(x)p_F(y)\pmod2.
\]

More generally, if $d=\dim M_F$ and $N_x(a)$ is the label multiplicity from
Theorem~\ref{thm:cellwise-holonomy-fourier},
then for all cells $x,y$,
\[
 \sum_{m\in M_F}u_{F,m}(x)u_{F,m}(y)
   =2^d\sum_{a\in\mathbb F_2^d}N_x(a)N_y(a).
\]
Consequently,
\[
 \sum_{m\in M_F\cap\{3,5,6\}}u_{F,m}(x)u_{F,m}(y)
 =2^d\sum_aN_x(a)N_y(a)-9t_F(x)t_F(y).             \tag{*}
\]

There is no gauge-selection rule of the preceding kind that is defined for
all Latin squares and is isotopy-equivariant, even when the complete
endpoint-lift data of Theorem~\ref{thm:endpoint-characterization} are
supplied.
\end{theorem}

\begin{proof}
A gauge change multiplies $u_{F,m}$ by one sign, which occurs twice in
$Q_{F,m}$.  The kernel assertion follows from
Theorem~\ref{thm:exact-mask-cell-pushforward}.  The two relevant families of
line indicators span a $(2n-1)$-dimensional space: each family has dimension
$n$, and their intersection consists only of the constant vector because the
two Latin partitions meet once in every pair of blocks.  This gives the rank
bound.  For even $m$, the role formula in
Theorem~\ref{thm:gauge-invariant-cell-lift} gives
$u_{F,m}(x)\equiv t_F(x)\pmod2$, proving the congruence.

For the Fourier identity, insert the character expansions at $x$ and $y$.
The sum over all characters of $\mathbb F_2^d$ is $2^d$ when the two labels
agree and zero otherwise.  Mask zero contributes
$u_{F,0}(x)u_{F,0}(y)=9t_F(x)t_F(y)$, while every odd-weight exact mask
vanishes by Theorem~\ref{thm:gauge-invariant-cell-lift}; subtracting the zero
mask gives~$(*)$.

For the final assertion, use the Klein-four table
\[
\begin{pmatrix}
0&1&2&3\\1&0&3&2\\2&3&0&1\\3&2&1&0
\end{pmatrix}.
\]
Fix the intercalate component on rows, columns, and symbols $\{0,1\}$.
The autotopism that fixes rows and applies $(0\,1)(2\,3)$ simultaneously to
columns and symbols stabilizes this component.  In a flag order for which it
acts as $(f_0\,f_1)(f_2\,f_3)$, normalized potentials for masks $3,5,6$ are
\[
 (0,1,1,0),\qquad(0,1,0,1),\qquad(0,0,1,1).
\]
The autotopism complements the first two potentials and fixes the third.
Adding a gauge constant does not change these shifts.  An equivariant
selection for mask $3$ or $5$ would therefore have to be both fixed and
complemented by the same autotopism, a contradiction.
Theorem~\ref{thm:endpoint-characterization} reconstructs the
endpoint lifts functorially from the table and conversely reconstructs the
table from them, so supplying those lifts cannot remove this symmetry.
\end{proof}

The theorem gives the canonical degree-two replacement for an impossible
equivariant linear lift.  It also exposes its exact boundary: components with
zero-dimensional exact-mask space have no nontrivial $u_{F,m}$ and remain
invisible to~$(*)$ even when their cross-component parity support from
Theorem~\ref{thm:cross-component-cell-parity} is nonzero.  Thus the
quadratic kernel is not an order-$10$ obstruction by itself.

\begin{theorem}[Component cell-edge partition]
\label{thm:component-cell-edge-partition}
For every connected flag component $F$ and every
$V\in\{R,C,S\}$ there is a simple graph $A_F^V$ on the Latin cells such
that
\[
 \bigsqcup_F E(A_F^V)=E(A^V),
\]
where $A^V$ is the disjoint union of the $n$ copies of $K_n$ on the
$V$-lines.  The three graphs belonging to one component have the same degree
vector:
\[
 \deg_{A_F^R}(x)=\deg_{A_F^C}(x)=\deg_{A_F^S}(x)=t_F(x).
 \tag{1}
\]
If $B_F$ is the flag-by-cell incidence matrix restricted to $F$ and
$A_F=A_F^R+A_F^C+A_F^S$, then
\[
 B_F^{\mathsf T}B_F=3\operatorname{diag}(t_F)+2A_F. \tag{2}
\]
Moreover, each
$\Lambda_F^V=\operatorname{diag}(t_F)-A_F^V$ is positive semidefinite and
\[
 \sum_F\Lambda_F^V=(n-1)I-A^V. \tag{3}
\]

On every $V$-line the entries of $t_F$ form the degree sequence of a simple
graph and hence satisfy all Erdos--Gallai inequalities.  The union $A_F$ is
connected on its nonisolated cell support.  For an intercalate component the
three $A_F^V$ are the three perfect matchings of $K_4$; consequently
\[
 \sum_{F\text{ nonintercalate}}\deg_{A_F^V}(x)=n-1-\iota(x) \tag{4}
\]
in every view.  Summing (4) over the views recovers the multiplicity-one
endpoint value
$((E_R\otimes E_C\otimes E_S)X)(x)=3(n-1-\iota(x))$ from
Theorem~\ref{thm:triple-endpoint-parity}.
\end{theorem}

\begin{proof}
Two flags joined across a dual edge of colour $V$ share exactly the two Latin
cells forming the corresponding $V$-edge.  Conversely, any two cells on one
$V$-line lie together in exactly two flags, and those two flags are paired by
that dual involution.  Dual edges therefore map bijectively to Latin-cell
edges.  Restricting this bijection to connected components gives the stated
edge partitions.

For a fixed cell $x$, the colour-$R$ involution pairs its occurrences in
roles $0$ and $2$, the colour-$C$ involution pairs roles $1$ and $2$, and the
colour-$S$ involution pairs roles $0$ and $1$.  Each role occurs $t_F(x)$
times by Theorem~\ref{thm:cell-role-affine-transport}.  Thus there are
$2t_F(x)$ incident flag sides of any fixed colour and two sides above every
cell edge, proving (1).

The diagonal entry of $B_F^{\mathsf T}B_F$ at $x$ counts the
$3t_F(x)$ flags of $F$ containing $x$.  Two distinct cells occur together in
two flags precisely when their shared cell edge belongs to $A_F$, and
otherwise never.  This proves (2).  Each $\Lambda_F^V$ is an ordinary graph
Laplacian.  Summing (1) over components gives
$\sum_F\operatorname{diag}(t_F)=(n-1)I$, while the edge partition gives
$\sum_FA_F^V=A^V$, proving (3).

All edges of $A_F^V$ stay within $V$-lines, so (1) restricted to one line is
an actual simple-graph degree sequence.  Its Erdos--Gallai inequalities are
therefore necessary; this is strictly stronger than even degree sum, since
$(n-1,n-1,0,\ldots,0)$ has even sum but violates the inequality at $k=2$
for $n\geq4$.  Every flag spans a triangle in $A_F$, and consecutive flags
along a connected dual path share a full cell edge, proving connectedness on
the support.

An intercalate component has four flags and four cells.  In every view (1)
gives degree one at each cell, hence a perfect matching; the three disjoint
matchings exhaust $K_4$.  Exactly $\iota(x)$ intercalate components contain
$x$.  Removing their unit contribution from the global degree $n-1$ proves
(4), and the final identity follows by summing the three views.
\end{proof}

This integral lift also applies to components whose exact-mask space has
dimension zero and therefore reaches the blind layer of
Theorem~\ref{thm:quadratic-gram-gauge-obstruction}.  It remains a necessary
projection: graphical line degrees alone do not ensure compatible rainbow
flag triangles or a common Latin-square realization.

\begin{theorem}[Orientable cell-companion surface]
\label{thm:cell-companion-surface}
Glue the rainbow Latin-cell triangle of every flag in a connected component
$F$ along the component cell edges of
Theorem~\ref{thm:component-cell-edge-partition}.  At a Latin cell $x$, its
link is a disjoint union of cycles.  Every link cycle has length divisible by
three, with edge-type vertices appearing periodically in the orders
$R,C,S$ or $R,S,C$.

Let $k_F(x)$ be the number of these link cycles.  Splitting $x$ into one
vertex for each link cycle gives a connected closed orientable surface
$\Sigma_F$ with
\[
 \chi(\Sigma_F)=\sum_x k_F(x)-\frac{|F|}{2}=2-2g_F. \tag{5}
\]
In particular,
\[
 \sum_xk_F(x)\equiv\frac{|F|}{2}\pmod2,
 \qquad
 \sum_xk_F(x)\leq\frac{|F|}{2}+2. \tag{6}
\]
For an intercalate component, $\Sigma_F$ is the sphere.
\end{theorem}

\begin{proof}
At $x$, take the incident component cell edges as link vertices.  Every flag
containing $x$ joins its two incident cell edges and hence supplies one link
edge.  A component cell edge belongs to two flags, so each link vertex has
degree two.  More precisely, an $R$-edge sees $x$ in roles $0$ and $2$ in
its two flags; its two link neighbours are therefore of types $S$ and $C$.
The cyclic statements hold for $C$- and $S$-edges.  Along any link cycle the
three types must consequently repeat in one of the two cyclic orders, proving
the divisibility assertion.

Orient each flag triangle by its cell-role order
$x_0\to x_1\to x_2\to x_0$.  Across $\alpha_S$ the common roles $0,1$ are
exchanged, across $\alpha_C$ the roles $1,2$ are exchanged, and across
$\alpha_R$ the roles $2,0$ are exchanged.  Adjacent faces therefore induce
opposite orientations on every common side.  Splitting vertices by link
components yields a closed orientable surface.  It remains connected because
the dual face-adjacency graph is the connected dual graph of $F$.

The normalized surface has $|F|$ faces, $3|F|/2$ edges, and
$\sum_xk_F(x)$ vertices, giving (5).  The parity and inequality in (6) follow
from integral nonnegative genus.  An intercalate has four faces and four
length-three cell links, so (5) gives genus zero.
\end{proof}

The cell-companion topology differs from the atom-vertex topology of
Theorem~\ref{thm:flag-surface}: an intercalate is a sphere in the former and
$\mathbb{RP}^2$ in the latter.  The new genus is a joint cross-view invariant,
but its necessary conditions still do not characterize Latin realizability.

\begin{theorem}[Normalized cell-incidence rank]
\label{thm:normalized-cell-rank}
Let $V_F$ be the normalized vertices of the cell-companion surface
$\Sigma_F$, and let $C_F$ be its flag-by-vertex incidence matrix.  Across
each glued triangle edge, join the two vertices opposite that edge.  If the
resulting graph $D_F$ on $V_F$ has $q_F$ connected components, then
\[
 \ker_{\mathbb Q}C_F
 =\{z:\text{$z$ is constant on the components of $D_F$ and one face
 sum is zero}\},                                      \tag{7}
\]
and hence
\[
 \operatorname{null}_{\mathbb Q}C_F=q_F-1,
 \qquad
 \operatorname{rank}_{\mathbb Q}C_F=|V_F|-q_F+1.       \tag{8}
\]
For every nonempty connected flag component, $1\leq q_F\leq3$; no FFF
assumption is needed for this bound.
If $d_F(v)$ is the vertex-link length and $M_F$ is the edge-multiplicity
adjacency matrix of $\Sigma_F$, then
\[
 C_F^{\mathsf T}C_F=\operatorname{diag}(d_F)+2M_F.      \tag{9}
\]
Finally, let $R_F$ forget the normalized sheet over each Latin cell.  For the
component flag-by-cell incidence matrix $B_F$ from
Theorem~\ref{thm:component-cell-edge-partition},
\[
 B_F=C_FR_F,
 \qquad
 B_F^{\mathsf T}B_F
 =R_F^{\mathsf T}(C_F^{\mathsf T}C_F)R_F.              \tag{10}
\]
\end{theorem}

\begin{proof}
The three corners of a flag lie over distinct Latin cells, and normalization
preserves the underlying cell, so every row of $C_F$ has three distinct
ones.  Subtract the incidence equations of two faces adjacent across one
edge.  Their two shared vertices cancel, leaving equality of the two vertices
opposite the edge.  A kernel vector is therefore constant on every component
of $D_F$.

Conversely, crossing a face edge replaces only its opposite vertex by a
vertex in the same $D_F$-component.  Thus a vector constant on those
components has the same three-vertex sum on every face.  The dual face graph
is connected, so all equations reduce to one face equation.  The resulting
linear functional on the $q_F$ component constants is nonzero over
$\mathbb Q$: the all-one vector has face sum $3$.  This proves (7) and (8).

More precisely, the multiset of $D_F$-component labels on the three corners
of a face is unchanged across any dual edge: the shared corners remain,
and the opposite corners are joined in $D_F$.  Dual connectedness propagates
this same multiset to every face.  Every vertex occurs in a face, so at most
three components can occur; nonemptiness gives $1\leq q_F\leq3$.

The diagonal Gram entry counts the faces at a normalized vertex and equals
its link length.  Two distinct vertices occur together in exactly the two
faces incident with each surface edge between them, proving (9).  Each
normalized vertex lies over one Latin cell; forgetting that sheet in every
face gives $B_F=C_FR_F$, and the Gram aggregation in (10) follows.
\end{proof}

Thus Theorem~\ref{thm:component-cell-edge-partition} is the coarse cell
aggregation of a canonical positive-semidefinite
surface Gram.  The exact rank need not be full: the frozen controls contain
components with $q_F=2$ and $q_F=3$.  Those finite distributions are not a
general classification, and the theorem remains a necessary componentwise
identity rather than an order-$10$ decision.

The atom and normalized-cell models also admit an exact mixed-rank
decomposition.

\begin{theorem}[Joint atom/cell incidence rank decomposition]
\label{thm:joint-atom-cell-rank}
For a connected flag component $F$, let $P_F$ be its flag-by-atom incidence,
let $C_F$ be its normalized flag-by-vertex incidence, and put
\[
 a_F=\#\{\text{atoms of }F\},\qquad
 v_F=|V_F|,\qquad
 s_F=\dim_{\mathbb Q}(\operatorname{im}P_F\cap
                      \operatorname{im}C_F).
\]
For $J_F=[P_F\ C_F]$ one has
\[
 \operatorname{rank}_{\mathbb Q}J_F
   =a_F+v_F-q_F-1-s_F,
 \qquad
 \operatorname{null}_{\mathbb Q}J_F=q_F+1+s_F.       \tag{11}
\]
Moreover $s_F\geq1$.  Thus, with $\epsilon_F=s_F-1$,
\[
 \operatorname{rank}_{\mathbb Q}J_F
   =a_F+v_F-q_F-2-\epsilon_F,
 \qquad
 \operatorname{null}_{\mathbb Q}J_F=q_F+2+\epsilon_F.\tag{12}
\]
The joint Gram $J_F^{\mathsf T}J_F$ is positive semidefinite and has the
same rank and nullity.  The integer $\epsilon_F$ is invariant under
relabelling flags, atoms, and normalized vertices.
\end{theorem}

\begin{proof}
Write $U_F=\operatorname{im}P_F$ and
$W_F=\operatorname{im}C_F$.  The exact atom projection on one connected
component gives $\dim U_F=a_F-2$, while
Theorem~\ref{thm:normalized-cell-rank} gives
$\dim W_F=v_F-q_F+1$.  Therefore
\[
 \operatorname{rank}J_F=\dim(U_F+W_F)
 =\dim U_F+\dim W_F-\dim(U_F\cap W_F),
\]
which is the rank formula in (11); subtracting it from $a_F+v_F$ gives the
nullity formula.

Every flag contains exactly one row atom, so assigning potential one to all
row atoms and zero to the other atoms gives the all-one face function in
$U_F$.  Every flag also has three normalized vertices, so assigning
potential $1/3$ to every normalized vertex gives the same all-one face
function in $W_F$.  Hence $s_F\geq1$, and (12) follows.  Finally
$\ker(J_F^{\mathsf T}J_F)=\ker J_F$ over $\mathbb Q$.  Relabellings only
permute rows or columns of the two incidence matrices, proving invariance.
\end{proof}

The exhaustive structural audit certifies the rational intersection
dimension for all $104645$ components in the frozen $9777$-table control
corpus.  In the $9768$ order-$8$ FFF components, the values
$s_F=1,2,3,4,13$ occur with multiplicities $9755,5,5,2,1$; the last case has
$\epsilon_F=12$.  The sorted component profile
$(|F|,a_F,v_F,q_F,s_F)$ takes $201$ values on the $230$ frozen order-$8$ FFF
representatives, with $176$ singleton values and maximum collision size
three.  Combining it with the canonical holonomy profile of
Computational Theorem~\ref{cthm:holonomy-profile} gives $202$ values and
$178$ singletons, still with maximum collision size three.  These are exact
finite-corpus invariants, not a main-class classification or a universal
bound on $s_F$.

The intersection in Theorem~\ref{thm:joint-atom-cell-rank} has a direct
cycle-space description.  For a flag $f$ and a colour $V$, write $a_V(f)$
for its $V$-atom, let $\alpha_V(f)$ be the adjacent flag across the edge of
colour $V$, and let $v_V(f)$ be the normalized vertex opposite that edge.
Define the atom-transition multigraph $\Gamma_F^V$ on the $V$-atoms by
joining $a_V(f)$ to $a_V(\alpha_V(f))$ for every dual $V$-edge.  Loops and
parallel edges are retained.

\begin{theorem}[Gradient characterization of common modes]
\label{thm:common-mode-gradient}
Each $\Gamma_F^V$ is connected.  For a normalized-vertex potential $z$, put
on an oriented dual $V$-edge $(f,g)$
\[
 \delta_z^V(f,g)=z(v_V(f))-z(v_V(g)).
\]
Then
\[
 C_Fz\in\operatorname{im}P_F
 \quad\Longleftrightarrow\quad
 \delta_z^R,\delta_z^C,\delta_z^S
 \text{ are exact gradients on }
 \Gamma_F^R,\Gamma_F^C,\Gamma_F^S.                 \tag{13}
\]
Equivalently, the signed sum of each cochain vanishes around every cycle of
its transition multigraph.

Let $X_F$ be the set of Latin cells used by $F$, let $B_F^X$ be the
flag-by-cell incidence restricted to these columns, and form the coarse
opposite-cell graph $E_F$ by joining the two underlying cells opposite each
dual edge.  If $p_F$ is the number of components of $E_F$, then
\[
 \operatorname{null}_{\mathbb Q}B_F^X=p_F-1,
 \qquad
 \operatorname{rank}_{\mathbb Q}B_F^X=|X_F|-p_F+1. \tag{14}
\]
The sheet-forgetting projection gives $1\leq p_F\leq q_F\leq3$ and
$\operatorname{im}B_F^X\subseteq\operatorname{im}C_F$.  Hence, with
\[
 t_F=\dim_{\mathbb Q}(\operatorname{im}P_F\cap
                       \operatorname{im}B_F^X),
 \qquad \eta_F=s_F-t_F,
\]
one has
\[
 1\leq t_F\leq s_F,
 \qquad
 0\leq\eta_F\leq |V_F|-|X_F|-q_F+p_F.             \tag{15}
\]
Thus $t_F-1$ is the nonconstant common-mode dimension already visible on
unsplit cells, while $\eta_F$ is the quotient dimension contributed only by
the normalization sheets.
\end{theorem}

\begin{proof}
Project a dual path in $F$ to its $V$-atoms.  Crossing a non-$V$ edge leaves
the $V$-atom fixed, while crossing a $V$-edge traverses an edge of
$\Gamma_F^V$; hence each transition graph is connected.

If $C_Fz=P_Fx$, subtract the equations of adjacent flags
$g=\alpha_V(f)$.  The two shared normalized vertices and the two shared
atoms cancel, leaving
\[
 z(v_V(f))-z(v_V(g))
   =x_V(a_V(f))-x_V(a_V(g)).
\]
Thus all three cochains are gradients.  Conversely choose atom potentials
$x_V$ realizing the three gradients.  The residual
\[
 h(f)=(C_Fz)(f)-\sum_Vx_V(a_V(f))
\]
is unchanged across every dual edge.  It is therefore constant because the
dual graph is connected; adding that constant to all atom potentials of one
colour makes the residual zero.  This proves (13), and the cycle criterion
is the standard characterization of exact cochains on a connected graph.

The same subtraction argument for $B_F^Xw=0$ says that $w$ is constant on
every component of $E_F$.  Conversely such a potential has the same
three-cell sum on every flag, so only one face-sum equation remains.  This
functional is nonzero over $\mathbb Q$, because the all-one potential has
sum three, proving (14).  Forgetting the sheet maps every edge of the
normalized opposite-vertex graph onto the corresponding edge of $E_F$, so
$p_F\leq q_F$; it also gives $B_F^X=C_FR_F$ after restricting to $X_F$.
The constant face mode gives $t_F\geq1$.  Finally, for
$U=\operatorname{im}P_F$, $T=\operatorname{im}B_F^X$, and
$W=\operatorname{im}C_F$, the natural map
$(U\cap W)/(U\cap T)\hookrightarrow W/T$ is injective.  Taking dimensions
and using (14) and Theorem~\ref{thm:normalized-cell-rank} proves (15).
\end{proof}

If $Z_F^{\rm norm}$ and $Z_F^{\rm cell}$ denote the normalized and coarse
potential spaces satisfying the corresponding three gradient conditions,
then the exact sequences induced by $C_F$ and $B_F^X$ give
\[
 s_F=\dim Z_F^{\rm norm}-(q_F-1),
 \qquad
 t_F=\dim Z_F^{\rm cell}-(p_F-1).                  \tag{16}
\]
There is also a rational Schur formulation.  For the orthogonal projectors
$\Pi_C,\Pi_B$ onto $\operatorname{im}C_F$ and
$\operatorname{im}B_F^X$, respectively,
\[
 \operatorname{null}\bigl(P_F^{\mathsf T}(I-\Pi_C)P_F\bigr)=2+s_F,
 \qquad
 \operatorname{null}\bigl(P_F^{\mathsf T}(I-\Pi_B)P_F\bigr)=2+t_F.
\]
This positive-semidefinite projector statement is asserted over
$\mathbb Q$ (or $\mathbb R$), not over arbitrary finite fields.

The complete frozen audit checks all $104645$ components of the $9777$
control tables and certifies $1111$ nonconstant integer coarse witnesses with
no error.  All $1080$ exceptional order-$6$ components have
$(s_F,t_F,\eta_F)=(2,2,0)$.  Among the $9768$ order-$8$ FFF components,
$t_F=1,2,3,4,13$ occurs with multiplicities $9756,5,4,2,1$.  Exactly one
observed component has positive normalization-only quotient: source $1955$,
component $4$, has $(s_F,t_F,\eta_F)=(3,1,2)$.  In particular, the source
$114$ component with $s_F=13$ is entirely coarse.  Every observed component
has $p_F=q_F$, but neither this equality nor uniqueness of the sheet-only
case is asserted universally.  The partial order-$10$ records remain
diagnostic only.

The componentwise gradient conditions have a particularly rigid global
intersection.

\begin{theorem}[Global common modes and exact Schur nullity]
\label{thm:global-common-mode}
Let the flag surface have $c$ connected components, let $a$ be its total
number of atoms, and let $P,B,Q$ be the global flag-by-atom,
flag-by-cell, and cell-by-atom support matrices.  Put
\[
 Z_L=\{w\in\mathbb Q^{n^2}:Bw\in\operatorname{im}P\},\qquad
 r_L=\dim Z_L,\qquad \kappa_L=r_L-1.
\]
For $n>3$, the map $w\mapsto Bw$ identifies $Z_L$ with
$\operatorname{im}P\cap\operatorname{im}B$, and
\[
 \operatorname{rank}_{\mathbb Q}[P\ B]=a-2c+n^2-r_L,
 \qquad
 \operatorname{null}_{\mathbb Q}[P\ B]=2c+r_L.       \tag{17}
\]
If $\Pi_B$ is the rational orthogonal projector onto
$\operatorname{im}B$, then
\[
 \operatorname{null}_{\mathbb Q}
 \bigl(P^{\mathsf T}(I-\Pi_B)P\bigr)
 =2c+r_L=2c+1+\kappa_L.                              \tag{18}
\]

Every common pair $P\xi=Bw$ has one common line sum $\lambda$ on every row,
column, and symbol class, and satisfies
\[
 Q\xi=(n-3)w+2\lambda\mathbf1.                       \tag{19}
\]
After setting
\[
 w_0=nw-\lambda\mathbf1,\qquad
 \xi_0=n\xi-3\lambda e_R,
\]
where $e_R$ is one on all row atoms and zero elsewhere, one has
\[
 P\xi_0=Bw_0,\qquad N^{\mathsf T}w_0=0,
 \qquad Q\xi_0=(n-3)w_0.                             \tag{20}
\]
Consequently $\kappa_L\leq(n-1)(n-2)$.  Finally, for the atom-only operator
\[
 \mathcal D_L=(n-3)P-BQ
\]
one has the exact formula
\[
 \operatorname{null}_{\mathbb Q}\mathcal D_L=2c+\kappa_L.    \tag{21}
\]
Thus the forced Schur lower bound $2c+1$ is exact precisely when the stacked
global gradient system has only the constant cell mode.
\end{theorem}

\begin{proof}
The Gram spectrum in Theorem~\ref{thm:flag-surface} makes $B$ injective for
$n>3$.  The image-sum formula, together with
$\operatorname{rank}P=a-2c$, then gives (17).  The map
$\ker(P^{\mathsf T}(I-\Pi_B)P)\to
\operatorname{im}P\cap\operatorname{im}B$, $x\mapsto Px$, has kernel
$\ker P$ and is onto, proving (18).

Write a cell as $x=(r,c,s)$.  Sum $P\xi=Bw$ over each of the three families
of $n-1$ flags in which $x$ occupies one fixed cell-corner role.  On the
atom side, every family enumerates once all atoms incident with $x$, and so
has sum $(Q\xi)_x$.  On the cell side the three sums are
\[
 (n-3)w_x+R_r+C_c,\quad
 (n-3)w_x+R_r+T_s,\quad
 (n-3)w_x+C_c+T_s,
\]
where $R,C,T$ are the row, column, and symbol line sums of $w$.  Equality
forces $R_r=C_c=T_s$ at every cell.  Since every row-column pair is a cell,
all $3n$ sums equal one value $\lambda$, and (19) follows.  The identities
$Pe_R=\mathbf1$ and $Qe_R=(n-1)\mathbf1$ now give (20).

The line-incidence Gram $N^{\mathsf T}N$ has diagonal blocks $nI$ and
off-diagonal blocks $J$, hence rank $3n-2$.  Therefore
$\dim\ker N^{\mathsf T}=n^2-(3n-2)=(n-1)(n-2)$, proving the bound.

Equation (20) maps centered common modes into $\ker\mathcal D_L$.  Conversely,
if $\mathcal D_L\xi=0$, put $w=Q\xi/(n-3)$.  Then $P\xi=Bw$; comparison with
(19) forces $\lambda=0$.  The resulting map from $\ker\mathcal D_L$ onto the
centered common-mode space has kernel $\ker P$, because
$Q=\tfrac13B^{\mathsf T}P$.  Its kernel has dimension $2c$ and its image
dimension $\kappa_L$, proving (21).
\end{proof}

The frozen global audit evaluates both descriptions of $r_L$ over two large
prime fields and supplies independent integer lower-bound witnesses.  All
$230$ order-$8$ FFF representatives have $r_L=1$, so their global Schur
nullity is exactly $2c+1$.  In the complete order-$6$ corpus,
$r_L=1,2,3$ occurs with multiplicities $8088,1080,240$.  All $1560$
nonconstant integer basis witnesses satisfy (19)--(20), and the atom-only
nullity (21) agrees over both prime fields on all $9777$ audited tables.
These distributions are finite-corpus statements; the theorem is general.
The $135$ tracked order-$10$ tables all have $r_L=1$, but that corpus is
partial and gives no order-$10$ decision.

The phrase ``trade space'' in Theorem~\ref{thm:global-common-mode} refers to
the integer dependencies among the occupied Latin triples.  It is useful to
separate this object from a classical Latin bitrade: a nonempty bitrade mate
changes symbols at fixed row--column positions and therefore uses ambient
triples outside the occupied set of the original square.

For a Latin square $L$, write
\[
 \mathcal T_L=\ker_{\mathbb Z}N^{\mathsf T}.
\]
If disjoint sets $A,B$ of four occupied Latin cells have the same incidence
with every row, column, and symbol line, call
$\mathbf1_A-\mathbf1_B$ a balanced $4+4$ cell exchange and let
$\mathcal E_4(L)$ be the lattice generated by all such exchanges.

\begin{computationaltheorem}[Order-$6$ common-mode exchange lattice]
\label{cthm:order6-common-mode-lattice}
Across all $9408$ reduced Latin squares of order $6$, the lattice
$\mathcal E_4(L)$ is saturated in its rational span.  It has rank $20$ for
$5448$ squares and rank $19$ for $3960$ squares.  Every balanced $3+3$ cell
exchange belongs to $\mathcal E_4(L)$.

For all $1080$ global-common-mode exceptions with $r_L=2$, and for $60$ of
the $240$
exceptions with $r_L=3$, one has $\mathcal E_4(L)=\mathcal T_L$.  In the
remaining $180$ cases,
\[
 \mathcal T_L/\mathcal E_4(L)\cong\mathbb Z,
\]
and the saturation in $\mathcal T_L$ of the full rational centred
common-mode space maps onto the index-two subgroup.  In each of these $180$
cases, a balanced $6+6$ cell exchange completes $\mathcal E_4(L)$ to
$\mathcal T_L$.
\end{computationaltheorem}

\begin{proof}
For each table, the program groups all three- and four-cell subsets by their
complete row/column/symbol incidence profile and pairs disjoint subsets in
the same group.  The line-incidence matrix has rank $16$, so
$\mathcal T_L$ has rank $20$.  Smith and Hermite normal forms in saturated
kernel coordinates determine the rank, saturation, containment, and quotient
indices.  The rational centred common-mode basis is saturated inside
$\mathcal T_L$ before the index-two comparison, so the conclusion is
independent of the particular integer certificates produced by the global
common-mode audit.

An independent implementation recomputes ambient Smith and Hermite forms
with Python-FLINT and checks every rank over $\mathbb F_2$, $\mathbb F_3$,
and $\mathbb F_{1000003}$.  It also verifies all $1560$ centred modes and
all $180$ stored $6+6$ completions.  Both implementations report zero
errors; the frozen commands, outputs, and hashes are listed in
Section~\ref{sec:reproducibility}.
\end{proof}

This theorem exhibits an integral layer that rational Schur rank does not
see.  It is a complete order-$6$ census result, not a universal quotient
theorem and not an order-$10$ obstruction.

Reduce $\mathcal E_4(L)$ modulo two and let $\Lambda_L$ be the binary span of
the row-, column-, and symbol-line indicators.  Since every balanced exchange
has even intersection with every line,
$\Lambda_L\subseteq\operatorname{Ann}(\overline{\mathcal E}_4)$.  If the
nonzero Smith factors of $N^{\mathsf T}$ are odd and $\mathcal E_4(L)$ is a
saturated codimension-one sublattice of $\mathcal T_L$, with rank preserved
modulo two, then


\[
 \operatorname{Ann}(overline{\mathcal E}_4)/\Lambda_L\cong\mathbb F_2.
 \tag{22}
\]

The unique nonzero class evaluates on $\mathcal T_L$ independently of its
representative, because line indicators pair trivially with line-balanced
vectors.  It is the unique binary character of
$\mathcal T_L/\mathcal E_4(L)$.  Isotopies and parastrophes transport cells,
lines, and balanced exchanges, so this class is relabelling covariant.  In
particular, the weight enumerator of its covector coset and the line-degree
profiles of minimum representatives are main-class invariants.

\begin{computationaltheorem}[Order-$6$ intrinsic exchange parity]
\label{cthm:order6-intrinsic-exchange-parity}
The complete order-$6$ census has $12$ main classes.  Exactly three, containing
$3960$ reduced squares, have $\operatorname{rank}\mathcal E_4(L)=19$.  Their
triples
\[
 (\text{reduced-class size},\text{minimum covector weight},
   \text{number of minima})
\]
are respectively
\[
 (3240,7,12),\qquad(540,9,80),\qquad(180,9,160).
\]
The last class is exactly the rank-$19$, $r_L=3$ family of Computational
Theorem~\ref{cthm:order6-common-mode-lattice}; all its representatives are
TTT and non-group-isotopic.  Every saturated centred common mode has parity
zero, while every stored $6+6$ completion has parity one.
\end{computationaltheorem}

\begin{proof}
Canonical reduction under all isotopies and six parastrophes partitions the
$9408$ reduced squares into class sizes
\[
20,40,60,108,180,360,540,1080,1080,1080,1620,3240.
\]
For every rank-$19$ table, exact binary elimination computes (22), enumerates
its $2^{16}$ covector coset, and records minimum weights and line profiles.
The three resulting signatures are constant inside, and distinct between,
the three classes.  An independent implementation rebuilds the balanced
$4+4$ supports directly and recomputes all three representative signatures;
it also checks the full source partition and reports zero errors.  The parity
statements for common modes and completions follow from the index-two and
primitive-completion conclusions of Computational
Theorem~\ref{cthm:order6-common-mode-lattice} and are checked on all $180$
tables.
\end{proof}

This identifies the index-two phenomenon as one precise order-$6$ main-class
structure.  It does not force such a quotient for an FFF square of order $10$.

For an arbitrary integer matrix $N$ with $m$ columns and an integral
sublattice $E\subseteq\ker_{\mathbb Z}N$, put
\[
 \rho_{\mathbb Q}=\operatorname{rank}_{\mathbb Q}N,\quad
 \rho_2=\operatorname{rank}_{\mathbb F_2}\overline N,\quad
 e_{\mathbb Q}=\operatorname{rank}_{\mathbb Q}E,\quad
 e_2=\dim_{\mathbb F_2}\overline E.
\]

\begin{proposition}[Binary exchange-defect decomposition]
\label{prop:exchange-defect-decomposition}
The quotient
\[
 Q_2(N,E)=\ker_{\mathbb F_2}\overline N/\overline E
\]
has dimension
\[
 \dim Q_2(N,E)
 =(m-\rho_{\mathbb Q}-e_{\mathbb Q})
  +(\rho_{\mathbb Q}-\rho_2)
  +(e_{\mathbb Q}-e_2).
 \tag{23}
\]
All three summands are nonnegative.  Moreover
\[
 \operatorname{Ann}(\overline E)/\operatorname{row}(\overline N)
 \cong Q_2(N,E)^*.
\]
\end{proposition}

\begin{proof}
Rank--nullity gives $\dim Q_2=m-\rho_2-e_2$; adding and subtracting
$\rho_{\mathbb Q}+e_{\mathbb Q}$ gives (23).  Reduction modulo a prime
cannot increase the rank of an integer matrix, and
$E\subseteq\ker_{\mathbb Q}N$, proving nonnegativity.  The annihilator
statement is the standard duality between a subspace of a nondegenerate
binary pairing and the dual of the corresponding quotient.
\end{proof}

For a Latin square, take $N=N_L$ and $E=\mathcal E_4(L)$, and denote the
dimension in (23) by $q_4(L)$.  Proposition
\ref{prop:exchange-defect-decomposition} separates three genuinely different
sources of a binary character: rational exchange codimension, line-incidence
rank loss modulo two, and exchange-rank loss modulo two.  These quantities
are main-class invariants, because isotopies and parastrophes merely permute
occupied cells and line rows.

\begin{computationaltheorem}[Order-$8$ FFF exchange-parity controls]
\label{cthm:order8-exchange-parity-controls}
Across the complete frozen set of $230$ order-$8$ FFF representatives,
\[
 \#\{L:q_4(L)=0,1,2,3\}=(39,147,35,9).
\]
The binary line ranks are distributed as
\[
 19^1\,20^{10}\,21^{156}\,22^{63},
\]
and the binary $\mathcal E_4$ ranks as
\[
 39^1\,40^8\,41^{46}\,42^{175}.
\]
In particular, FFF does not force a nonzero or a codimension-one binary
exchange quotient.  One explicit non-group-isotopic FFF witness is frozen
local source line $170$ (ANU zero-based record $169$), for which
\[
 \operatorname{rank}_{\mathbb F_2}\overline N_L=22,\qquad
 \dim_{\mathbb F_2}\overline{\mathcal E}_4(L)=42,\qquad q_4(L)=0.
\]
\end{computationaltheorem}

\begin{proof}
For each table, the audit enumerates all four-cell subsets, groups them by
their complete row/column/symbol incidence profile, and pairs exactly the
disjoint subsets with equal profile.  Direct binary elimination computes the
two ranks and verifies both inclusions defining the quotient.  A separate
packed-profile implementation rebuilds all exchanges and ranks for all
$230$ tables and reports zero discrepancies.  It also stores the full
local-source-line-$170$ table and independently confirms its FFF census
identity.
\end{proof}

On the separate, nonexhaustive corpus of $135$ tracked order-$10$ diagnostic
tables, every record has $\rho_2=28$, $e_2=72$, and hence $q_4=0$; four
pattern representatives were independently recomputed.  Those records are
not FFF examples and do not contribute to the order-$10$ nonexistence proof.
Their role is diagnostic:
the exchange-parity layer is already saturated on the best available
near-models, so it supplies no additional obstruction there.

\begin{definition}[Fibres, contractions, and coherent edge lifts]
For a finite set $A$, write $E(A)$ for its unordered two-subsets.  Given
$T:E(A)\times E(B)\times E(C)\to\{0,1,4\}$, a two-label fibre is obtained by
fixing two edge labels and varying the third; its type is the multiset of its
positive values.  For $a\in A$ and $f\in E(B)$, the $A$-endpoint contraction is
\[
 C^a(f,g)=\sum_{\substack{e\in E(A)\\a\in e}}T(e,f,g)
 \qquad(g\in E(C)),
\]
with cyclic analogues in the other two colours.  Its dominant output is the
unique $g$ at which the contraction has value at least two.  The dominant
outputs for a fixed endpoint $a$ are a coherent two-subset lift if there is a
point permutation $\lambda_a:B\to C$ such that the dominant output for every
$f\in E(B)$ is $\lambda_a(f)=\{\lambda_a(b):b\in f\}$.  The corresponding
families in the other colours are denoted $\mu_b:A\to C$ and
$\nu_c:A\to B$.
\end{definition}

\begin{theorem}[Endpoint characterization]
\label{thm:endpoint-characterization}
Let $A,B,C$ be sets of cardinality $n\geq3$.  A function
\[
 T:E(A)\times E(B)\times E(C)\longrightarrow\{0,1,4\}
\]
is the pair-label tensor of a Latin square if and only if the following hold:
every two-label fibre is empty, $[1,1]$, or $[4]$; every endpoint contraction
in all three colours is $[2,1,1]$ or $[4]$; and, for each fixed endpoint, its
dominant-output map on all input edges is coherently the two-subset lift of one
point permutation.  The three recovered permutation families are
automatically sharply transitive and define the same Latin square.  Adding
the link-divisibility condition of
Proposition~\ref{prop:pair-label-tensor} characterizes the tensors of FFF
Latin squares.
\end{theorem}

\begin{proof}
Necessity is Proposition~\ref{prop:pair-label-tensor}.  For sufficiency, let
$X$ and $Y$ denote the tensor entries of multiplicity one and four.  Fix a
$B$-edge $f$.  If $p_f$ fibres $T(*,f,g)$ have type $[1,1]$ and $q_f$ have
type $[4]$, double counting their incidences with endpoints of $A$ gives
$p_f+2q_f=n$.  There are also exactly $n$ dominant endpoint contractions.
A $[4]$ fibre contributes two dominant endpoints, whereas the two distinct
$A$-edges in a $[1,1]$ fibre contribute at most their one common endpoint.
Equality forces those two edges to meet.  Cyclically, every $x\in X$ therefore
has unique dominant endpoints $a(x),b(x),c(x)$ in its three edge labels.

Let $\lambda_a:B\to C$ and $\mu_b:A\to C$ be the point permutations supplied
by the first two endpoint families.  For fixed $(a,b)$ count
\[
 D_{ab}=\#\{x\in X:a(x)=a, b(x)=b\}
 +\#\{(e,f,g)\in Y:a\in e, b\in f\}.
\]
Every $X$-entry is counted once and every $Y$-entry four times, so
$\sum_{a,b}D_{ab}=n^2(n-1)$.  An object counted by $D_{ab}$ has a $C$-edge
$g$ containing both $\lambda_a(b)$ and $\mu_b(a)$.  The map to $g$ is
injective, since $(a,b,g)$ recovers $f$ and $e$ through the two bijective
edge lifts.  Two vertex stars of $K_n$ intersect in $n-1$ edges when their
centres agree and in one edge otherwise.  Hence $D_{ab}\leq n-1$ for all
$a,b$.  The total attains the sum of these $n^2$ upper bounds, and $n-1>1$;
therefore
\[
 \lambda_a(b)=\mu_b(a)\qquad(a\in A,b\in B).
\]
The cyclic count for $(a,c)$ similarly gives
$\nu_c(a)=\lambda_a^{-1}(c)$.  Thus
$L(a,b)=\lambda_a(b)=\mu_b(a)$ has permutational rows and columns and is
Latin; all three families are compatible and sharply transitive.

Finally, a multiplicity-one entry with dominant corner $(a,b,c)$ has
$c\ne L(a,b)$: equality would make the corresponding rectangle an
intercalate and force the second edge in a $[1,1]$ fibre to equal the first.
Each multiplicity-four entry supplies its four nonincidence corners.  Indeed,
both endpoint rows map its column edge bijectively onto its symbol edge, and
column uniqueness makes the resulting rectangle an intercalate.  Conversely,
a multiplicity-one entry cannot reconstruct an intercalate: its second row
endpoint would then also dominate the same symbol edge, so the two distinct
support edges in the relevant $[1,1]$ fibre would contain both row endpoints
and hence coincide.  The preceding star injection is a bijection, so the
resulting corners exhaust exactly the $n^2(n-1)$ triples $(a,b,c)$ with
$c\ne L(a,b)$.  The recovery map
\[
 (a,b,c)\longmapsto
 \bigl(\{a,\mu_b^{-1}(c)\},\{b,\lambda_a^{-1}(c)\},
       \{L(a,b),c\}\bigr)
\]
returns each selected corner to its original three pair labels by the
dominant edge-lift equations.  The flag bijection of
Theorem~\ref{thm:flag-surface} and the preceding
dichotomy therefore give one flag over each entry of $X$, four over each
entry of $Y$, and none elsewhere.  Hence $T$ is exactly the flag tensor of
$L$.  The FFF statement now follows from the link criterion in
Proposition~\ref{prop:pair-label-tensor}.
\end{proof}

The restriction $n\geq3$ is sharp for faithful recovery of the witnessing
point permutations and their automatic compatibility.  At $n=2$ every edge
set is a singleton, so the two-subset action cannot distinguish the two point
permutations.  The sole tensor value four is nevertheless the flag tensor of
the unique order-$2$ Latin square.

\begin{proposition}[Intrinsic recognition of endpoint edge lifts]
\label{prop:intrinsic-edge-lift}
Let $U,W$ be $n$-element sets and let $M:E(U)\to E(W)$ be a bijection.  If
$n=3$ or $n\geq5$, then $M$ is the two-subset action of a unique point
bijection $U\to W$ if and only if it preserves edge intersection.  If $n=4$,
the intersection-preserving maps are the point-induced maps and their
compositions with target-edge complementation $e\mapsto W\setminus e$;
requiring vertex stars to map to vertex stars selects exactly the
point-induced maps.

Consequently the point-lift axiom in
Theorem~\ref{thm:endpoint-characterization} may be replaced, for $n=3$ or
$n\geq5$, by bijectivity and preservation of intersection on pair labels.
At $n=4$ one adds preservation of the family of vertex stars.
\end{proposition}

\begin{proof}
In the line graph of a complete graph, cliques are pairwise-intersecting edge
families.
Such a family either has a common endpoint or has at most three members:
after choosing $\{a,b\}$ and $\{a,c\}$, any member avoiding $a$ must be
$\{b,c\}$.  Hence, for $n\geq5$, the cliques of size $n-1$ are precisely the
vertex stars.  An intersection-preserving edge bijection maps the source
stars onto the target stars, defining a point bijection $p:U\to W$.  Since
$\{\{u,v\}\}=\operatorname{Star}(u)\cap\operatorname{Star}(v)$, it follows
that $M(\{u,v\})=\{p(u),p(v)\}$.  Uniqueness is immediate.  At $n=3$, there
are six point bijections and six edge bijections, and the two-subset action is
faithful.

For $n=4$, edge complementation preserves intersection but sends each vertex
star to the complementary triangle, so it is not point-induced.  The three
complementary edge pairs are exactly the nonedges of $L(K_4)$.  Every graph
automorphism permutes these three pairs and may swap within them, yielding at
most $3!2^3=48$ maps.  The 24 point lifts and their 24 distinct complement
compositions attain this bound.  Exactly the first coset preserves the
vertex stars.
\end{proof}

The order-$4$ exception is also a sharp warning against replacing endpoint
recognition by local completion rules alone.  Complementing one pair-label
coordinate of a genuine order-$4$ tensor preserves its fibre trichotomy and
all coordinatewise edge-intersection relations, but exhaustive comparison
with all $576$ labelled order-$4$ Latin squares computes that the transformed
tensor is not Latin-realizable.

\begin{definition}[Cell-local cycle-incidence signature]
\label{def:cell-local-cycle-signature}
For a Latin cell $x=(r,c,L(r,c))$, a view
$V\in\{R,C,S\}$, and $2\leq\ell\leq n$, let $a^V_\ell(x)$ count the other
lines whose induced permutation with the distinguished $V$-line through $x$
has an $\ell$-cycle containing $x$.  Put
\[
 A^V(x)=(a^V_2(x),\ldots,a^V_n(x)).
\]
The \emph{cell-local cycle-incidence signature} is the multiset
\[
 \bigl\{(A^R(x),A^C(x),A^S(x)):x\text{ a Latin cell}\bigr\},
\]
canonically minimized over the six simultaneous permutations of the three
view coordinates.  This convention makes the signature invariant under
isotopy and parastrophy.  Pointwise
$\sum_\ell a^V_\ell(x)=n-1$, and the three values
$a^R_2(x),a^C_2(x),a^S_2(x)$ all equal the number of intercalates through
$x$.
\end{definition}

\begin{computationaltheorem}[Order-$8$ rainbow-square classifier]
\label{cthm:rainbow-classifier}
On the $230$ order-$8$ FFF main-class representatives, the pair consisting of
the signature in Definition~\ref{def:cell-local-cycle-signature} and
\[
 w_6(L)=\operatorname{Fix}((\alpha_R\alpha_C\alpha_S)^2)
\]
takes $230$ distinct values.
\end{computationaltheorem}

The cell-local signature alone has $226$ values and four collision pairs.
Their respective $w_6$ values are
$(352,288)$, $(448,192)$, $(448,192)$, and $(288,352)$.  The complete
flag identities were audited with no error on all $9408$ reduced order-$6$
squares and all $283657$ order-$8$ representatives.  A separate
implementation recomputed $w_6$ on all $230$ FFF representatives with zero
mismatches.  This is a finite census classifier, not an injectivity theorem
for higher orders.

\begin{computationaltheorem}[Pair-label triangle classifier]
\label{cthm:pair-triangle-classifier}
On the same $230$ representatives, the pair consisting of the
cell-local cycle-incidence signature from
Definition~\ref{def:cell-local-cycle-signature} and $t(L)$ takes $230$
distinct values.
\end{computationaltheorem}

The scalar $t(L)$ alone takes $100$ values and $(I(L),t(L))$ takes $188$.
The four cell-local-signature collision pairs $[4,22]$, $[11,114]$,
$[13,19]$, and
$[117,118]$ have respective $t$-values
$(2112,1856)$, $(2560,1536)$, $(2048,1536)$, and $(1792,1920)$.
The pair matrices were reconstructed both from flag codegrees and from the
induced edge lifts, with no mismatch.  As above, this is an exact classifier
only on the frozen order-$8$ FFF census.
